
\section{Preguntas de reflexión}

\subsubsection*{¿Qué ventajas ofrece RMI sobre sockets tradicionales?}

RMI permite que el código cliente invoque métodos definidos en una interfaz Java aunque el objeto se encuentre en otro proceso o equipo. El programador trabaja con operaciones como `listarNoticias()` o `editarNoticia(...)`, en vez de diseñar manualmente un formato de mensajes, delimitadores, códigos de operación y reglas para leer bytes parciales desde un socket.

El stub y el runtime abstraen el intercambio de mensajes. RMI serializa argumentos, resultados y excepciones; el Registry ofrece un mecanismo de localización; y `RemoteException` hace visible que una llamada que parece ordinaria puede fallar por la red. Las interfaces remotas también constituyen un contrato verificable por el compilador: cliente y servidor conocen las firmas y los tipos compartidos.

Con sockets tradicionales existe mayor libertad sobre el protocolo, la representación binaria, el uso de memoria, las conexiones y la compatibilidad entre lenguajes. Esa libertad es valiosa para protocolos especializados o de alto rendimiento, pero obliga a implementar framing, serialización, despacho, correlación de respuestas, evolución del protocolo y manejo de errores. RMI reduce ese trabajo para un sistema Java a Java y hace la práctica más cercana al dominio.

La abstracción no elimina la naturaleza distribuida: una invocación puede tardar, fallar después de haber ejecutado parcialmente una acción o devolver una excepción. Además, RMI acopla ambos extremos a Java y a clases compatibles, su interoperabilidad es menor, la configuración de puertos/firewall requiere atención y la serialización debe tratarse como una frontera sensible. Por ello RMI es más productivo que sockets para este ejercicio homogéneo, pero no siempre es el transporte preferido para sistemas abiertos o políglotas.

\subsubsection*{¿Cómo se manejaron múltiples peticiones simultáneas?}

El runtime de RMI puede atender distintas invocaciones concurrentemente. La implementación no supone que los métodos se ejecuten uno tras otro ni sincroniza globalmente el objeto remoto. En su lugar, las clases de servicio y repositorio son esencialmente sin estado: los parámetros y variables de cada petición permanecen locales al hilo que la atiende.

Cada operación obtiene una conexión JDBC propia, y cada escritura delimita su transacción. De esa manera una solicitud no comparte un `Connection`, `PreparedStatement` o `ResultSet` mutable con otra. H2 coordina el acceso a sus datos y las transacciones evitan estados parciales. Las lecturas se mantienen independientes y no esperan un monitor global de Java.

Las sesiones sí son estado compartido, por lo que se almacenan en un `ConcurrentHashMap`. La creación, consulta, expiración y eliminación de tokens se realizan con operaciones seguras de esa colección. Las noticias no se duplican en un mapa del servidor: H2 continúa siendo su única fuente de verdad.

Las pruebas de concurrencia previstas usan `ExecutorService`, `CountDownLatch` y futuros para iniciar lectores, publicaciones y ediciones con suficiente simultaneidad. En particular, dos tareas que editan la misma versión deben producir exactamente un éxito y un conflicto, sin depender de pausas arbitrarias con `Thread.sleep`.

\subsubsection*{¿Qué problemas enfrentaron al implementar consistencia de los datos?}

El caso principal es la actualización perdida. Dos redactores pueden abrir la misma noticia, conservar la misma versión y enviar cambios distintos. Si el servidor ejecutara un `UPDATE` solo por identificador, el último reemplazaría al primero sin saber que trabajó sobre información obsoleta. El campo `version` y una condición adicional en el `WHERE` convierten esa carrera en un conflicto explícito.

También se deben coordinar lecturas mientras existen escrituras. Las transacciones hacen que una publicación o edición sea una unidad: no se expone una noticia con solo parte de sus campos modificados. El orden por fecha de modificación e identificador produce resultados repetibles incluso cuando dos filas comparten una marca temporal.

La consistencia no se limita a concurrencia. El servidor debe comprobar que el autor de la sesión sea propietario de la noticia, pues confiar en un botón deshabilitado permitiría modificar datos mediante otro cliente. Las claves foráneas impiden autores inexistentes; `NOT NULL`, usuario único y `CHECK (version >= 1)` refuerzan invariantes; la validación limita textos y rechaza entradas vacías; y las consultas preparadas separan datos de SQL.

Finalmente, una excepción a mitad de una escritura debe provocar rollback. La dificultad está en traducir el fallo a una excepción de dominio comprensible sin ocultarlo en logs ni filtrar detalles internos. La solución separa el registro técnico del mensaje remoto y utiliza `try-with-resources` para liberar recursos aun cuando la operación falla.

\subsubsection*{¿Cómo escalarías este sistema en un entorno real?}

Primero separaría la lógica de negocio del transporte para conservarla al reemplazar RMI por gRPC o una API REST según los consumidores, requisitos de interoperabilidad y operación. Para clientes internos fuertemente tipados, gRPC permitiría contratos eficientes y multiplataforma; para integración web y ecosistemas amplios, HTTP/REST podría ser más accesible. La elección no sería automática: dependería de latencia, streaming, compatibilidad y herramientas disponibles.

H2 se sustituiría por una base administrada con pool de conexiones, migraciones, respaldos, réplicas y estrategia de recuperación. Varias instancias del servicio se colocarían detrás de un balanceador. Las sesiones dejarían de residir en un único proceso: podrían usarse tokens verificables de corta duración o un almacén compartido y altamente disponible. Las escrituras requerirían claves de idempotencia para soportar reintentos ambiguos.

Para cargas de lectura elevadas podrían incorporarse cachés con invalidación explícita y réplicas de lectura. La publicación de eventos mediante mensajería desacoplaría indexación, auditoría y notificaciones. Un motor de búsqueda atendería consultas de texto completo, y WebSocket, Server-Sent Events o streaming proporcionarían notificaciones en tiempo real.

La operación necesitaría observabilidad: logs estructurados y centralizados, métricas, trazas distribuidas, tableros y alertas. La seguridad incluiría TLS mutuo o una terminación TLS confiable, proveedor de identidad, rotación de secretos, mínimos privilegios, auditoría y límites de solicitudes. Réplicas en zonas distintas, comprobaciones de salud, despliegues automatizados y pruebas de recuperación aportarían alta disponibilidad.

Ninguna de estas mejoras se implementa en la práctica porque añadiría infraestructura que ocultaría el objetivo: demostrar Java RMI, serialización, persistencia local, concurrencia y fallos con la arquitectura mínima que permite estudiarlos.
